\documentclass[12pt,a4paper]{article}

\usepackage[a4paper,top=2.2cm,bottom=2.2cm,left=2.5cm,right=2.5cm]{geometry}
\usepackage{fontspec}
\usepackage{tikz}
\usetikzlibrary{arrows.meta,calc}
\usepackage{enumitem}
\usepackage{setspace}
\usepackage{hyperref}
\usepackage{xcolor}
\usepackage{fancyhdr}

\setmainfont[
  Path=fonts/,
  UprightFont=Sarabun-Regular.ttf,
  BoldFont=Sarabun-Bold.ttf,
  ItalicFont=Sarabun-Italic.ttf,
  BoldItalicFont=Sarabun-BoldItalic.ttf,
  Ligatures=TeX
]{Sarabun}

\XeTeXlinebreaklocale "th"
\XeTeXlinebreakskip=0pt plus 1pt minus 0.5pt
\hyphenpenalty=10000
\exhyphenpenalty=10000
\emergencystretch=3em
\hypersetup{colorlinks=true,linkcolor=black,urlcolor=black,citecolor=black}
\urlstyle{same}
\setlength{\parindent}{1.25cm}
\setlength{\parskip}{0pt}
\setlist[enumerate]{nosep,leftmargin=1.15cm,labelsep=0.35cm}
\setlist[itemize]{nosep,leftmargin=1.15cm}
\pagestyle{fancy}
\fancyhf{}
\renewcommand{\headrulewidth}{0pt}
\fancyfoot[C]{\fontsize{12}{14}\selectfont \thepage}

\newcommand{\mainsection}[1]{%
  \par\vspace{0.65\baselineskip}%
  \noindent\textbf{#1}\par\nobreak\vspace{0.2\baselineskip}%
}

\newcommand{\subsectionthai}[1]{%
  \par\vspace{0.35\baselineskip}%
  \noindent\textbf{#1}\par\nobreak\vspace{0.1\baselineskip}%
}

\newcommand{\schedulearrow}[3]{%
  \pgfmathsetmacro{\arrowy}{-1.40-(#1-0.5)*1.22}%
  \pgfmathsetmacro{\arrowxstart}{8.08+(#2-1)*0.66+0.12}%
  \pgfmathsetmacro{\arrowxend}{8.08+#3*0.66-0.12}%
  \draw[{Latex[length=2.0mm]}-{Latex[length=2.0mm]},line width=0.55pt]
    (\arrowxstart,\arrowy) -- (\arrowxend,\arrowy);
}

\begin{document}
\fontsize{16}{20}\selectfont
\onehalfspacing

\begin{center}
  {\fontsize{18}{22}\selectfont\bfseries โครงร่างโครงงานวิจัย}\par
  \vspace{0.65\baselineskip}
  {\fontsize{17}{21}\selectfont\bfseries
  การศึกษาทำซ้ำและประเมินแบบจำลองการจำแนกระยะการนอนหลับจากสัญญาณคลื่นไฟฟ้าสมองช่องสัญญาณเดียวโดยใช้การแปลงซินโครสควีซและการเรียนรู้เชิงลึกแบบลำดับ}
\end{center}

\mainsection{1. ที่มาและความสำคัญของโครงงาน}

การนอนหลับเป็นกระบวนการทางสรีรวิทยาที่มีความสำคัญต่อการฟื้นฟูร่างกาย การทำงานของระบบประสาท การเรียนรู้ ความจำ และการควบคุมอารมณ์ การเปลี่ยนแปลงของร่างกายในระหว่างการนอนหลับสามารถจำแนกออกเป็นระยะต่าง ๆ ได้แก่ ระยะตื่น ระยะหลับไม่ฝัน และระยะหลับฝัน แต่ละระยะมีลักษณะของสัญญาณทางสรีรวิทยาที่แตกต่างกัน การจำแนกระยะการนอนหลับอย่างถูกต้องจึงมีความสำคัญต่อการศึกษาคุณภาพการนอนหลับและความผิดปกติที่เกี่ยวข้องกับการนอนหลับ

โดยทั่วไป การตรวจประเมินการนอนหลับดำเนินการด้วยวิธีพอลิซอมโนกราฟี (polysomnography: PSG) ซึ่งเป็นการบันทึกสัญญาณทางสรีรวิทยาหลายชนิดพร้อมกัน เช่น สัญญาณคลื่นไฟฟ้าสมอง สัญญาณการกลอกตา และสัญญาณไฟฟ้ากล้ามเนื้อ จากนั้นผู้เชี่ยวชาญจะตรวจสอบสัญญาณและกำหนดระยะการนอนหลับของช่วงสัญญาณแต่ละช่วงด้วยตนเอง อย่างไรก็ตาม กระบวนการดังกล่าวต้องอาศัยผู้เชี่ยวชาญ ใช้เวลานาน และอาจเกิดความแตกต่างในการให้คะแนนระหว่างผู้ประเมินได้

สัญญาณคลื่นไฟฟ้าสมอง (electroencephalogram: EEG) เป็นสัญญาณที่แสดงการเปลี่ยนแปลงของกิจกรรมทางไฟฟ้าของสมอง และมีรูปแบบแตกต่างกันตามระยะการนอนหลับ ด้วยเหตุนี้ จึงมีการนำเทคนิคการประมวลผลสัญญาณและการเรียนรู้ของเครื่องมาประยุกต์ใช้เพื่อจำแนกระยะการนอนหลับโดยอัตโนมัติ โดยเฉพาะการใช้สัญญาณ EEG เพียงช่องสัญญาณเดียว ซึ่งช่วยลดจำนวนอิเล็กโทรดและความซับซ้อนของระบบตรวจวัด

Aram และคณะ [1] ได้นำเสนอแบบจำลองการจำแนกระยะการนอนหลับจากสัญญาณ EEG ช่องสัญญาณ Fpz-Cz โดยแบ่งสัญญาณออกเป็นช่วงละ 30 วินาที และแปลงสัญญาณจากโดเมนเวลาไปเป็นภาพแทนความสัมพันธ์ระหว่างเวลาและความถี่ด้วยการแปลงซินโครสควีซ (synchrosqueezed transform: SST) จากนั้นจึงเรียนรู้คุณลักษณะของสัญญาณด้วยโครงข่ายประสาทเทียมแบบคอนโวลูชัน (convolutional neural network: CNN) ร่วมกับการเรียนรู้แบบมีผู้กำกับเชิงเปรียบเทียบ (supervised contrastive learning) และใช้โครงข่ายหน่วยความจำแบบเกตเวียนกลับ (gated recurrent unit: GRU) เพื่อเรียนรู้ความสัมพันธ์ระหว่างระยะการนอนหลับที่เกิดต่อเนื่องกัน ผลการศึกษารายงานความถูกต้องในการจำแนกประมาณร้อยละ 83 [1]

แม้ว่างานวิจัยดังกล่าวจะแสดงให้เห็นถึงประสิทธิภาพของการผสานการวิเคราะห์สัญญาณในโดเมนเวลา-ความถี่กับการเรียนรู้เชิงลึก แต่การนำวิธีการดังกล่าวมาใช้ซ้ำยังมีประเด็นที่ควรตรวจสอบ ทั้งด้านขั้นตอนการเตรียมข้อมูล การกำหนดพารามิเตอร์ของแบบจำลอง และวิธีแบ่งข้อมูลสำหรับการฝึกและการประเมินผล โดยเฉพาะการแบ่งข้อมูลในระดับช่วงสัญญาณ ซึ่งอาจทำให้ข้อมูลจากบุคคลเดียวกันปรากฏอยู่ทั้งในชุดฝึกและชุดทดสอบ ส่งผลให้ค่าประสิทธิภาพที่ได้อาจยังไม่สะท้อนความสามารถของแบบจำลองในการจำแนกข้อมูลจากบุคคลที่ไม่เคยปรากฏในชุดฝึกอย่างแท้จริง

ดังนั้น การศึกษาครั้งนี้จึงมุ่งทำซ้ำกระบวนการจำแนกระยะการนอนหลับตามแนวทางของ Aram และคณะ โดยใช้ชุดข้อมูล Sleep-EDF Database Expanded จาก PhysioNet [2], [3] พร้อมทั้งประเมินแบบจำลองภายใต้แนวทางการแบ่งข้อมูลสองรูปแบบ ได้แก่ การแบ่งข้อมูลในระดับช่วงสัญญาณตามงานวิจัยต้นฉบับ และการแบ่งข้อมูลโดยแยกผู้เข้าร่วมการศึกษาออกจากกันอย่างชัดเจน การดำเนินการดังกล่าวจะช่วยตรวจสอบความสามารถในการทำซ้ำของวิธีการ ตลอดจนประเมินความสามารถของแบบจำลองในการประยุกต์ใช้กับข้อมูลจากบุคคลรายใหม่ได้อย่างเหมาะสมยิ่งขึ้น

\mainsection{2. วัตถุประสงค์ของโครงงาน}

\subsectionthai{2.1}
เพื่อศึกษาหลักการและสร้างกระบวนการเตรียมข้อมูลสัญญาณคลื่นไฟฟ้าสมองสำหรับการจำแนกระยะการนอนหลับตามแนวทางของงานวิจัยต้นฉบับ

\subsectionthai{2.2}
เพื่อพัฒนาแบบจำลองการจำแนกระยะการนอนหลับจากสัญญาณ EEG ช่องสัญญาณเดียว โดยใช้การแปลงซินโครสควีซ โครงข่ายประสาทเทียมแบบคอนโวลูชัน การเรียนรู้แบบมีผู้กำกับเชิงเปรียบเทียบ และโครงข่ายหน่วยความจำแบบ GRU

\subsectionthai{2.3}
เพื่อประเมินความสามารถในการทำซ้ำผลการศึกษาของงานวิจัยต้นฉบับภายใต้การแบ่งข้อมูลในระดับช่วงสัญญาณ

\subsectionthai{2.4}
เพื่อประเมินความสามารถของแบบจำลองในการจำแนกระยะการนอนหลับของบุคคลที่ไม่เคยปรากฏอยู่ในชุดข้อมูลสำหรับฝึกแบบจำลอง

\subsectionthai{2.5}
เพื่อเปรียบเทียบผลการประเมินระหว่างการแบ่งข้อมูลในระดับช่วงสัญญาณกับการแบ่งข้อมูลในระดับบุคคล

\subsectionthai{2.6}
เพื่อวิเคราะห์ประสิทธิภาพของแบบจำลองในแต่ละระยะการนอนหลับ รวมทั้งศึกษาข้อจำกัดและปัจจัยที่ทำให้ผลการศึกษาทำซ้ำแตกต่างจากผลที่รายงานในงานวิจัยต้นฉบับ

\mainsection{3. ขอบเขตของโครงงาน}

\subsectionthai{3.1 ขอบเขตด้านชุดข้อมูล}

ใช้ข้อมูลจากชุดข้อมูล Sleep-EDF Database Expanded รุ่น 1.0.0 ซึ่งเผยแพร่โดย PhysioNet โดยเลือกข้อมูลส่วน Sleep Cassette ประกอบด้วยการบันทึกการนอนหลับจำนวน 153 ระเบียน จากผู้เข้าร่วมการศึกษาจำนวน 78 คน [2], [3]

ข้อมูลจากผู้เข้าร่วมจำนวน 6 คนที่เคยใช้ในการทดลองนำร่อง จะใช้สำหรับตรวจสอบความถูกต้องของโปรแกรมและความพร้อมของระบบเท่านั้น และจะไม่นำผลจากข้อมูลดังกล่าวมาใช้เป็นข้อสรุปแทนผลการประเมินจากชุดข้อมูลฉบับสมบูรณ์

\subsectionthai{3.2 ขอบเขตด้านสัญญาณ}

ใช้สัญญาณ EEG ช่องสัญญาณ Fpz-Cz ซึ่งบันทึกด้วยอัตราการสุ่มตัวอย่าง 100 เฮิรตซ์ แบ่งสัญญาณออกเป็นช่วงละ 30 วินาที หรือ 3,000 จุดข้อมูลต่อช่วง

จำแนกผลออกเป็น 5 กลุ่ม ได้แก่

\begin{enumerate}[label=\arabic*.]
  \item ระยะตื่น (Wake: W)
  \item ระยะหลับไม่ฝันระยะที่ 1 (N1)
  \item ระยะหลับไม่ฝันระยะที่ 2 (N2)
  \item ระยะหลับไม่ฝันระยะที่ 3 (N3)
  \item ระยะหลับฝัน (rapid eye movement: REM)
\end{enumerate}

ข้อมูลระยะหลับเดิม S3 และ S4 จะถูกรวมเป็นระยะ N3 ส่วนข้อมูลที่กำกับว่าเป็นช่วงการเคลื่อนไหวหรือช่วงที่ไม่สามารถให้คะแนนได้จะไม่นำมาใช้ในการฝึกแบบจำลอง

\subsectionthai{3.3 ขอบเขตด้านการเตรียมข้อมูล}

ดำเนินการตรวจสอบความครบถ้วนของแฟ้มสัญญาณและแฟ้มคำกำกับระยะการนอนหลับ จับคู่ข้อมูลแต่ละระเบียน และตรวจสอบหมายเลขผู้เข้าร่วมการศึกษา ก่อนนำสัญญาณมาแบ่งเป็นช่วงเวลาและกำหนดฉลาก

ช่วงตื่นที่อยู่ห่างจากช่วงการนอนหลับมากเกินไปจะถูกตัดออกตามเกณฑ์ของงานวิจัยต้นฉบับ เพื่อลดความไม่สมดุลของจำนวนตัวอย่างระหว่างแต่ละระยะการนอนหลับ

\subsectionthai{3.4 ขอบเขตด้านการแปลงสัญญาณและแบบจำลอง}

สัญญาณ EEG แต่ละช่วงจะถูกแบ่งออกเป็นส่วนย่อย และแปลงเป็นภาพแทนสัญญาณในโดเมนเวลา-ความถี่ด้วยการแปลงซินโครสควีซ โดยใช้ส่วนจริงและส่วนจินตภาพของผลการแปลงเป็นข้อมูลนำเข้าสองช่องสัญญาณ

แบบจำลองประกอบด้วยกระบวนการเรียนรู้สามขั้น ดังนี้

\begin{enumerate}[label=\arabic*.]
  \item การฝึกโครงข่าย CNN เพื่อเรียนรู้คุณลักษณะพื้นฐานของระยะการนอนหลับ
  \item การปรับปรุงตัวแทนคุณลักษณะด้วยการเรียนรู้แบบมีผู้กำกับเชิงเปรียบเทียบ
  \item การนำตัวเข้ารหัสคุณลักษณะมาใช้ร่วมกับโครงข่าย GRU เพื่อเรียนรู้ความสัมพันธ์ตามลำดับเวลาของระยะการนอนหลับ
\end{enumerate}

\subsectionthai{3.5 ขอบเขตด้านการแบ่งข้อมูล}

การศึกษากำหนดแนวทางการประเมินแยกออกเป็นสองส่วนอย่างชัดเจน ได้แก่

\begin{enumerate}[label=\arabic*.]
  \item \textbf{การศึกษาทำซ้ำตามข้อความในงานวิจัยต้นฉบับ} แบ่งข้อมูลในระดับช่วงสัญญาณเป็นชุดฝึก ชุดตรวจสอบ และชุดทดสอบ เพื่อใช้ตรวจสอบว่ากระบวนการที่พัฒนาขึ้นสามารถให้ผลใกล้เคียงกับงานวิจัยต้นฉบับหรือไม่
  \item \textbf{การประเมินความสามารถในการจำแนกข้อมูลของบุคคลรายใหม่} แบ่งข้อมูลในระดับบุคคล โดยไม่อนุญาตให้ข้อมูลจากบุคคลเดียวกันปรากฏอยู่มากกว่าหนึ่งชุดข้อมูล เพื่อป้องกันการรั่วไหลของข้อมูลระหว่างชุดฝึกและชุดทดสอบ
\end{enumerate}

ผลจากการแบ่งข้อมูลทั้งสองรูปแบบจะรายงานแยกจากกัน และจะไม่นำค่าประสิทธิภาพมาเปรียบเทียบโดยตรงโดยไม่อธิบายความแตกต่างของวิธีแบ่งข้อมูล

\subsectionthai{3.6 ขอบเขตด้านการประเมินผล}

ประเมินประสิทธิภาพของแบบจำลองด้วยค่าความถูกต้อง ความแม่นยำ ความไว ค่าคะแนน F1 แบบเฉลี่ยมหภาคและเฉลี่ยถ่วงน้ำหนัก ค่าสัมประสิทธิ์แคปปาของโคเฮน และเมทริกซ์สับสน

นอกจากนี้ จะรายงานค่าคะแนน F1 ของระยะ W, N1, N2, N3 และ REM แยกจากกัน เนื่องจากจำนวนข้อมูลและความยากในการจำแนกของแต่ละระยะมีความแตกต่างกัน

\subsectionthai{3.7 ขอบเขตการนำผลไปใช้}

แบบจำลองที่พัฒนาขึ้นเป็นต้นแบบสำหรับการศึกษาและการประเมินเชิงวิจัยเท่านั้น ยังไม่ใช่อุปกรณ์ทางการแพทย์ และไม่สามารถนำไปใช้วินิจฉัยโรคหรือทดแทนการพิจารณาของแพทย์และผู้เชี่ยวชาญด้านการนอนหลับได้

\mainsection{4. วิธีดำเนินการวิจัย}

\subsectionthai{4.1 รูปแบบการวิจัย}

การศึกษาครั้งนี้เป็นการวิจัยเชิงทดลองด้านคอมพิวเตอร์ โดยใช้ข้อมูลทุติยภูมิที่เปิดเผยต่อสาธารณะจากฐานข้อมูล Sleep-EDF Database Expanded ไม่มีการรับสมัครผู้เข้าร่วมการศึกษาเพิ่มเติม และไม่มีการเก็บสัญญาณทางสรีรวิทยาจากบุคคลโดยตรง การดำเนินงานแบ่งเป็นการเตรียมข้อมูล การสร้างตัวแทนสัญญาณ การพัฒนาแบบจำลอง การฝึกแบบจำลอง การประเมินผล และการตรวจสอบความสามารถในการทำซ้ำ

\subsectionthai{4.2 การรวบรวมและตรวจสอบชุดข้อมูล}

ดาวน์โหลดข้อมูลส่วน Sleep Cassette จาก PhysioNet พร้อมตรวจสอบความครบถ้วนของแฟ้ม PSG และแฟ้ม Hypnogram ตามรายการข้อมูลต้นฉบับ ตรวจสอบรหัสผู้เข้าร่วม คืนที่บันทึก อัตราการสุ่มตัวอย่าง และค่าแฮชของแฟ้ม เพื่อให้สามารถตรวจสอบย้อนกลับได้ ชุดข้อมูลฉบับสมบูรณ์ต้องประกอบด้วย 153 ระเบียน จากผู้เข้าร่วม 78 คน ก่อนเริ่มการสร้างข้อมูล SST หรือการฝึกแบบจำลอง

\subsectionthai{4.3 การกำหนดระเบียบวิธีและการแบ่งข้อมูลล่วงหน้า}

กำหนดรายละเอียดของระเบียบวิธีไว้ล่วงหน้าก่อนการฝึกแบบจำลอง ได้แก่ วิธีจับคู่แฟ้ม วิธีจัดลำดับระเบียน การแมปฉลาก การตัดช่วงตื่น วิธีแบ่งข้อมูล ค่าเมล็ดสุ่ม และตัวชี้วัดที่ใช้ประเมินผล การประเมินแบ่งออกเป็น 2 ระเบียบวิธีหลัก คือ ระเบียบวิธีตามข้อความในงานวิจัยต้นฉบับและระเบียบวิธีที่แยกผู้เข้าร่วมระหว่างชุดฝึก ชุดตรวจสอบ และชุดทดสอบอย่างสมบูรณ์ หากรายละเอียดในบทความและโปรแกรมของผู้วิจัยต้นฉบับแตกต่างกัน จะบันทึกเป็นข้อแตกต่างของระเบียบวิธีและไม่นำผลจากคนละเงื่อนไขมารวมกัน

\subsectionthai{4.4 การเตรียมข้อมูลสัญญาณ}

เลือกสัญญาณ EEG ช่อง Fpz-Cz ที่อัตราการสุ่มตัวอย่าง 100 เฮิรตซ์ และแบ่งสัญญาณเป็นช่วงละ 30 วินาที กำหนดฉลากเป็น W, N1, N2, N3 และ REM โดยรวมระยะ S3 และ S4 เป็น N3 และตัดช่วง Movement หรือช่วงที่ไม่สามารถให้คะแนนได้ออก จากนั้นตัดช่วงตื่นที่อยู่นอกช่วงก่อนและหลังการนอนหลับตามเกณฑ์ที่กำหนด ตรวจสอบจำนวนช่วงสัญญาณและจำนวนตัวอย่างของแต่ละชั้นให้สอดคล้องกับหลักฐานอ้างอิงก่อนดำเนินการขั้นถัดไป

\subsectionthai{4.5 การสร้างภาพแทนสัญญาณด้วยการแปลงซินโครสควีซ}

แบ่งสัญญาณแต่ละช่วงออกเป็นหน้าต่างย่อยตามเงื่อนไขของระเบียบวิธี แล้วคำนวณการแปลงซินโครสควีซเพื่อสร้างภาพแทนสัญญาณในโดเมนเวลา-ความถี่ ใช้ส่วนจริงและส่วนจินตภาพเป็นข้อมูลนำเข้าสองช่องสัญญาณ ปรับขนาดและทำให้ค่าอยู่ในช่วงที่กำหนดเหมือนกันทุกชุดข้อมูล พร้อมตรวจสอบรูปร่างข้อมูล ค่าต่ำสุด ค่าสูงสุด ค่าที่ไม่เป็นจำนวน และค่าตรวจสอบของตัวอย่างมาตรฐานก่อนสร้างข้อมูลทั้งชุด

\subsectionthai{4.6 การพัฒนาและฝึกแบบจำลอง}

การฝึกแบบจำลองแบ่งเป็น 3 ขั้น ได้แก่

\begin{enumerate}[label=\arabic*.]
  \item \textbf{ขั้นที่ 1 การเรียนรู้คุณลักษณะด้วย CNN} ฝึกโครงข่าย CNN เพื่อจำแนกภาพแทนสัญญาณออกเป็น 5 ระยะการนอนหลับ และบันทึกตัวเข้ารหัสคุณลักษณะที่ผ่านการฝึกแล้ว
  \item \textbf{ขั้นที่ 2 การเรียนรู้แบบมีผู้กำกับเชิงเปรียบเทียบ} นำตัวเข้ารหัสจากขั้นที่ 1 มาปรับปรุงให้ตัวอย่างที่มีฉลากเดียวกันอยู่ใกล้กัน และตัวอย่างต่างฉลากอยู่ห่างกันในปริภูมิคุณลักษณะ
  \item \textbf{ขั้นที่ 3 การเรียนรู้ความสัมพันธ์ตามลำดับเวลา} นำตัวเข้ารหัสคุณลักษณะมาใช้ร่วมกับโครงข่าย GRU แบบซ้อน เพื่อเรียนรู้ความสัมพันธ์ระหว่างช่วงการนอนหลับที่ต่อเนื่องกันและทำนายฉลากของแต่ละช่วง
\end{enumerate}

ใช้จำนวนรอบการฝึก อัตราการเรียนรู้ ตัวปรับค่า และตัวกำหนดอัตราการเรียนรู้ตามระเบียบวิธีที่ลงทะเบียนไว้ สำหรับขั้นที่ 3 ใช้ขนาดชุดข้อมูลเชิงผลเท่ากับ 128 ตัวอย่าง หากข้อจำกัดของหน่วยความจำทำให้ต้องแบ่งเป็นชุดย่อย จะเฉลี่ยเกรเดียนต์จากตัวอย่างครบ 128 ตัวอย่างก่อนจำกัดบรรทัดฐานรวมของเกรเดียนต์ไว้ที่ 1.0 วิธีดังกล่าวจะระบุเป็นมาตรการเพิ่มเสถียรภาพของการคำนวณ และรายงานแยกจากเงื่อนไขของบทความต้นฉบับอย่างชัดเจน

\subsectionthai{4.7 การประเมินผลแบบจำลอง}

ประเมินแบบจำลองภายใต้การแบ่งข้อมูลในระดับช่วงสัญญาณและระดับบุคคล โดยใช้ค่าเมล็ดสุ่มอย่างน้อย 2 ค่า ได้แก่ 42 และ 43 เพื่อตรวจสอบความแปรปรวนจากการสุ่ม รายงานค่าความถูกต้อง ความแม่นยำ ความไว คะแนน F1 แบบเฉลี่ยมหภาคและเฉลี่ยถ่วงน้ำหนัก ค่าสัมประสิทธิ์แคปปาของโคเฮน คะแนน F1 รายระยะ และเมทริกซ์สับสน ผลจากการแบ่งข้อมูลคนละรูปแบบจะนำเสนอแยกตารางและไม่ตีความเป็นการเปรียบเทียบโดยตรง

\subsectionthai{4.8 การควบคุมคุณภาพและการตรวจสอบความทำซ้ำได้}

บันทึกรุ่นโปรแกรม ไลบรารี พารามิเตอร์ ค่าเมล็ดสุ่ม โครงสร้างแบบจำลอง ค่าแฮชของข้อมูล และแฟ้มผลลัพธ์ของการทดลองทุกครั้ง ตรวจสอบว่าไม่มีข้อมูลจากบุคคลเดียวกันรั่วไหลข้ามชุดข้อมูล ตรวจสอบค่าที่ไม่เป็นจำนวนและความสมบูรณ์ของจุดตรวจแบบจำลอง เมื่อการฝึกเสร็จสิ้นจะเรียกกระบวนการตรวจสอบผลแบบแยกจากกระบวนการฝึก และตรวจทานหลักฐานการทดลองก่อนนำผลไปสรุป

\subsectionthai{4.9 การวิเคราะห์และรายงานผล}

เปรียบเทียบผลที่ได้กับค่าที่รายงานในงานวิจัยต้นฉบับโดยพิจารณาความแตกต่างของชุดข้อมูล ขั้นตอนการเตรียมข้อมูล การแบ่งข้อมูล และพารามิเตอร์การฝึก วิเคราะห์สาเหตุของความคลาดเคลื่อนและประสิทธิภาพรายระยะ โดยให้ความสำคัญกับระยะ N1 ซึ่งมีจำนวนตัวอย่างน้อยและมีรูปแบบใกล้เคียงกับระยะข้างเคียง สรุปข้อจำกัดของการศึกษาและขอบเขตที่สามารถอ้างอิงผลได้โดยไม่กล่าวอ้างเกินกว่าหลักฐาน

\subsectionthai{4.10 ระยะเวลาดำเนินการ}

กำหนดระยะเวลาดำเนินงานรวม 12 เดือน โดยแต่ละกิจกรรมสามารถดำเนินการเหลื่อมกันได้ตามความพร้อมของข้อมูลและทรัพยากรคอมพิวเตอร์ ดังตารางที่ 1

\begin{center}
{\bfseries ตารางที่ 1 แผนการดำเนินงานตลอดระยะเวลา 12 เดือน}\par
\vspace{0.4\baselineskip}
\begin{tikzpicture}[x=1cm,y=1cm,scale=0.995,transform shape]
  \def\tablewidth{16.00}
  \def\tablebottom{-13.60}

  \draw[line width=0.55pt] (0,0) rectangle (\tablewidth,\tablebottom);
  \draw[line width=0.45pt] (8.08,0) -- (8.08,\tablebottom);
  \draw[line width=0.45pt] (8.08,-0.70) -- (\tablewidth,-0.70);
  \foreach \n in {1,...,11} {
    \pgfmathsetmacro{\vx}{8.08+\n*0.66}
    \draw[line width=0.35pt] (\vx,-0.70) -- (\vx,\tablebottom);
  }
  \foreach \r in {0,...,10} {
    \pgfmathsetmacro{\hy}{-1.40-\r*1.22}
    \draw[line width=0.35pt] (0,\hy) -- (\tablewidth,\hy);
  }

  \node[align=center,font=\fontsize{13}{15}\selectfont] at (4.04,-0.70)
    {กิจกรรม/ขั้นตอนการดำเนินงาน};
  \node[align=center,font=\fontsize{13}{15}\selectfont] at (12.04,-0.35) {เดือน};

  \foreach \m in {1,...,12} {
    \pgfmathsetmacro{\mx}{8.08+(\m-0.5)*0.66}
    \node[font=\fontsize{11.5}{13}\selectfont] at (\mx,-1.05) {\m};
  }

  \node[anchor=west,text width=7.63cm,align=left,font=\fontsize{11.3}{13.2}\selectfont]
    at (0.16,-2.01) {1. ศึกษาเอกสารที่เกี่ยวข้องและกำหนดระเบียบวิธีการวิจัย};
  \node[anchor=west,text width=7.63cm,align=left,font=\fontsize{11.3}{13.2}\selectfont]
    at (0.16,-3.23) {2. รวบรวม ตรวจสอบ และจัดทำบัญชีข้อมูล Sleep-EDF ฉบับสมบูรณ์};
  \node[anchor=west,text width=7.63cm,align=left,font=\fontsize{11.3}{13.2}\selectfont]
    at (0.16,-4.45) {3. เตรียมสัญญาณ กำหนดฉลาก และตรวจสอบจำนวนข้อมูล};
  \node[anchor=west,text width=7.63cm,align=left,font=\fontsize{11.3}{13.2}\selectfont]
    at (0.16,-5.67) {4. สร้างและตรวจสอบภาพแทนสัญญาณด้วยการแปลงซินโครสควีซ};
  \node[anchor=west,text width=7.63cm,align=left,font=\fontsize{11.3}{13.2}\selectfont]
    at (0.16,-6.89) {5. พัฒนาและฝึกแบบจำลอง CNN ขั้นที่ 1};
  \node[anchor=west,text width=7.63cm,align=left,font=\fontsize{11.3}{13.2}\selectfont]
    at (0.16,-8.11) {6. ฝึกการเรียนรู้แบบมีผู้กำกับเชิงเปรียบเทียบ ขั้นที่ 2};
  \node[anchor=west,text width=7.63cm,align=left,font=\fontsize{11.3}{13.2}\selectfont]
    at (0.16,-9.33) {7. พัฒนาและฝึกแบบจำลอง CNN-GRU ขั้นที่ 3};
  \node[anchor=west,text width=7.63cm,align=left,font=\fontsize{11.3}{13.2}\selectfont]
    at (0.16,-10.55) {8. ประเมินผลแบบช่วงสัญญาณและแบบแยกผู้เข้าร่วม};
  \node[anchor=west,text width=7.63cm,align=left,font=\fontsize{11.3}{13.2}\selectfont]
    at (0.16,-11.77) {9. วิเคราะห์ผล ตรวจสอบความทำซ้ำได้ และทบทวนข้อจำกัด};
  \node[anchor=west,text width=7.63cm,align=left,font=\fontsize{11.3}{13.2}\selectfont]
    at (0.16,-12.99) {10. สรุปผล จัดทำรายงานวิจัย และเตรียมต้นฉบับเผยแพร่};

  \schedulearrow{1}{1}{2}
  \schedulearrow{2}{1}{3}
  \schedulearrow{3}{2}{4}
  \schedulearrow{4}{3}{5}
  \schedulearrow{5}{4}{6}
  \schedulearrow{6}{5}{7}
  \schedulearrow{7}{6}{8}
  \schedulearrow{8}{7}{10}
  \schedulearrow{9}{9}{11}
  \schedulearrow{10}{11}{12}
\end{tikzpicture}
\end{center}

\mainsection{5. ประโยชน์ที่คาดว่าจะได้รับ}

\subsectionthai{5.1}
ได้กระบวนการเตรียมข้อมูลและแบบจำลองสำหรับจำแนกระยะการนอนหลับจากสัญญาณ EEG ช่องสัญญาณเดียวที่สามารถตรวจสอบย้อนกลับและทำซ้ำได้

\subsectionthai{5.2}
ได้ผลการประเมินความสามารถในการทำซ้ำของวิธีการจำแนกระยะการนอนหลับตามงานวิจัยต้นฉบับภายใต้ชุดข้อมูลและเงื่อนไขที่กำหนด

\subsectionthai{5.3}
ได้ข้อมูลเกี่ยวกับความแตกต่างระหว่างผลการประเมินด้วยการแบ่งข้อมูลในระดับช่วงสัญญาณกับการแบ่งข้อมูลในระดับบุคคล

\subsectionthai{5.4}
ได้แนวทางป้องกันการรั่วไหลของข้อมูลระหว่างชุดฝึกและชุดทดสอบสำหรับงานวิจัยที่ใช้ข้อมูลสัญญาณจากบุคคลหลายราย

\subsectionthai{5.5}
ได้ข้อมูลเกี่ยวกับข้อจำกัดของแบบจำลอง โดยเฉพาะการจำแนกระยะ N1 และระยะอื่นที่มีจำนวนตัวอย่างน้อยหรือมีรูปแบบสัญญาณใกล้เคียงกัน

\subsectionthai{5.6}
สามารถใช้ผลการศึกษาเป็นพื้นฐานสำหรับพัฒนาระบบวิเคราะห์การนอนหลับที่ใช้จำนวนช่องสัญญาณลดลง และสำหรับศึกษาการประยุกต์ใช้แบบจำลองกับข้อมูลจากบุคคลรายใหม่ในอนาคต

\mainsection{6. เอกสารอ้างอิง}

\begingroup
\fontsize{12.3}{14.5}\selectfont
\raggedright
\setlength{\parindent}{0pt}
\begin{enumerate}[label={[\arabic*]},leftmargin=1.15cm,itemsep=0.25\baselineskip,align=left]
  \item S. Aram, M. M. Ghassemi, and B. Mohammadzadeh Asl, ``Single-channel EEG sleep stage classification using synchrosqueezed transform and sequential representation learning,'' \textit{Scientific Reports}, 2026. \href{https://doi.org/10.1038/s41598-026-62832-5}{doi:10.1038/s41598-026-62832-5}
  \item B. Kemp, A. H. Zwinderman, B. Tuk, H. A. C. Kamphuisen, and J. J. L. Obery\'e, ``Analysis of a sleep-dependent neuronal feedback loop: The slow-wave microcontinuity of the EEG,'' \textit{IEEE Transactions on Biomedical Engineering}, vol. 47, no. 9, pp. 1185-1194, 2000. \href{https://doi.org/10.1109/10.867928}{doi:10.1109/10.867928}
  \item B. Kemp, ``Sleep-EDF Database Expanded,'' PhysioNet, version 1.0.0, 2013. \href{https://doi.org/10.13026/C2X676}{doi:10.13026/C2X676}
  \item P. Khosla et al., ``Supervised contrastive learning,'' in \textit{Advances in Neural Information Processing Systems}, vol. 33, pp. 18661-18673, 2020. \href{https://proceedings.neurips.cc/paper/2020/hash/d89a66c7c80a29b1bdbab0f2a1a94af8-Abstract.html}{เอกสารเผยแพร่}
  \item A. Supratak, H. Dong, C. Wu, and Y. Guo, ``DeepSleepNet: A model for automatic sleep stage scoring based on raw single-channel EEG,'' \textit{IEEE Transactions on Neural Systems and Rehabilitation Engineering}, vol. 25, no. 11, pp. 1998-2008, 2017. \href{https://doi.org/10.1109/TNSRE.2017.2721116}{doi:10.1109/TNSRE.2017.2721116}
\end{enumerate}
\endgroup

\end{document}
